% GENERATED FILE. Edit canonical lessons, then run npm run generate:books.
% canonical-source-hash: fnv1a64:dd38b3aefc1381c4
% canonical-lessons: ML-C532-laksanam, ML-C532-anubadha, ML-C532-sastrakriya, ML-C532-roganirnnayam, ML-C532-cikitsikkuka, ML-R532-review-525-532

\chapter{Symptoms and Surgery}
\label{ch:ml-symptoms-and-surgery}
\hlchaptermodality{532}

\begin{chapteropening}
\textbf{By the end of this chapter:} \emph{I can say a symptom, an infection, surgery, a diagnosis and to treat.}

Complete the last lesson of chapter 532: I can say a symptom, an infection, surgery, a diagnosis and to treat.
\end{chapteropening}

\section[\texorpdfstring{\ml{ലക്ഷണം} (lakṣaṇaṁ)}{ലക്ഷണം (lakṣaṇaṁ)}]{\ml{ലക്ഷണം} (lakṣaṇaṁ) — a symptom; a sign}

\label{lesson:ML-C532-laksanam}



\begin{quote}
Before the new one: say the Malayalam for showing off, then the Malayalam for to bother.
\end{quote}

\subsection*{You'll want to know: \ml{ലക്ഷണം}}
\textbf{\ml{ലക്ഷണം}} — \emph{lakṣaṇaṁ} — \textquotedblleft{}a symptom; a sign\textquotedblright{}.

\begin{quote}
\ml{പനിയും} \ml{ചുമയും} \ml{ആ} \ml{രോഗത്തിന്റെ} \ml{ലക്ഷണങ്ങളാണ്}.
\end{quote}

\begin{quote}
\emph{paniyuṁ cumayuṁ ā rōgattinṟe lakṣaṇaṅṅaḷāṇŭ.}
\end{quote}

\begin{quote}
Fever and a cough are symptoms of that illness.
\end{quote}

Plural: \ml{ലക്ഷണങ്ങൾ}, \emph{lakṣaṇaṅṅaḷ}.

It also means a sign or an omen: \ml{നല്ല} \ml{ലക്ഷണം}, a good sign.

\begin{cousinweb}
\textbf{\ml{ലക്ഷണം}} (\emph{lakṣaṇaṁ}) is borrowed from Sanskrit \emph{lakṣaṇa}, \textquotedblleft{}a line, a mark, a sign\textquotedblright{} (CDIAL 10882). A symptom is the mark an illness leaves.
\end{cousinweb}

\begin{grammarlens}[title={What you've built}]
One more word for describing what you lived and felt.
\end{grammarlens}

\subsection*{Your turn}
\begin{itemize}
\raggedright
  \item \emph{Say it:} \emph{lakṣaṇaṁ}
  \item \emph{Say it:} \emph{lakṣaṇaṁ}, once more
  \item \emph{Say it:} \emph{paniyuṁ cumayuṁ ā rōgattinṟe lakṣaṇaṅṅaḷāṇŭ.} — Fever and a cough are symptoms of that illness.
  \item \emph{Recall:} say the Malayalam for showing off, then the Malayalam for to bother, then say \emph{lakṣaṇaṁ} again
\end{itemize}

\begin{tcolorbox}[breakable,colback=teal!4,colframe=teal!35!black,title={Before you move on}]
What does \ml{ലക്ഷണം} mean? (\textquotedblleft{}A symptom; a sign\textquotedblright{}.) Say it once more.
\end{tcolorbox}
\section[\texorpdfstring{\ml{അണുബാധ} (aṇubādha)}{അണുബാധ (aṇubādha)}]{\ml{അണുബാധ} (aṇubādha) — an infection}

\label{lesson:ML-C532-anubadha}



\begin{quote}
Before the new one: say the Malayalam for to bother, then the Malayalam for a symptom.
\end{quote}

\subsection*{You'll want to know: \ml{അണുബാധ}}
\textbf{\ml{അണുബാധ}} — \emph{aṇubādha} — \textquotedblleft{}an infection\textquotedblright{}.

\begin{quote}
\ml{മുറിവിൽ} \ml{അണുബാധ} \ml{ഉണ്ടായി}.
\end{quote}

\begin{quote}
\emph{muṟivil aṇubādha uṇṭāyi.}
\end{quote}

\begin{quote}
The wound got infected.
\end{quote}

\ml{അണു} a germ + \ml{ബാധ} an affliction. \ml{അണു} also means an atom.

\begin{cousinweb}
\textbf{\ml{അണുബാധ}} (\emph{aṇubādha}) is a learned compound built from Sanskrit \emph{aṇu}, \textquotedblleft{}fine, minute\textquotedblright{} (CDIAL 193), and \emph{bādha}, \textquotedblleft{}affliction, distress\textquotedblright{} (Monier-Williams).
\end{cousinweb}

\begin{grammarlens}[title={What you've built}]
One more word for describing what you lived and felt.
\end{grammarlens}

\subsection*{Your turn}
\begin{itemize}
\raggedright
  \item \emph{Say it:} \emph{aṇubādha}
  \item \emph{Say it:} \emph{aṇubādha}, once more
  \item \emph{Say it:} \emph{muṟivil aṇubādha uṇṭāyi.} — The wound got infected.
  \item \emph{Recall:} say the Malayalam for to bother, then the Malayalam for a symptom, then say \emph{aṇubādha} again
\end{itemize}

\begin{tcolorbox}[breakable,colback=teal!4,colframe=teal!35!black,title={Before you move on}]
What does \ml{അണുബാധ} mean? (\textquotedblleft{}An infection\textquotedblright{}.) Say it once more.
\end{tcolorbox}
\section[\texorpdfstring{\ml{ശസ്ത്രക്രിയ} (śastrakriya)}{ശസ്ത്രക്രിയ (śastrakriya)}]{\ml{ശസ്ത്രക്രിയ} (śastrakriya) — surgery; an operation}

\label{lesson:ML-C532-sastrakriya}



\begin{quote}
Before the new one: say the Malayalam for a symptom, then the Malayalam for an infection.
\end{quote}

\subsection*{You'll want to know: \ml{ശസ്ത്രക്രിയ}}
\textbf{\ml{ശസ്ത്രക്രിയ}} — \emph{śastrakriya} — \textquotedblleft{}surgery; an operation\textquotedblright{}.

\begin{quote}
\ml{അച്ഛന്} \ml{അടുത്ത} \ml{ആഴ്ച} \ml{ഒരു} \ml{ശസ്ത്രക്രിയ} \ml{ഉണ്ട്}.
\end{quote}

\begin{quote}
\emph{acchanŭ aṭutta āḻca oru śastrakriya uṇṭŭ.}
\end{quote}

\begin{quote}
Father has an operation next week.
\end{quote}

The formal word; in speech many say \ml{ഓപ്പറേഷൻ}.

\begin{cousinweb}
\textbf{\ml{ശസ്ത്രക്രിയ}} (\emph{śastrakriya}) is a learned compound built from Sanskrit \emph{śastra}, \textquotedblleft{}an instrument for cutting\textquotedblright{} (CDIAL 12367), and \emph{kriyā}, \textquotedblleft{}work, an act\textquotedblright{} (CDIAL 3587): work with the knife.
\end{cousinweb}

\begin{grammarlens}[title={What you've built}]
One more word for describing what you lived and felt.
\end{grammarlens}

\subsection*{Your turn}
\begin{itemize}
\raggedright
  \item \emph{Say it:} \emph{śastrakriya}
  \item \emph{Say it:} \emph{śastrakriya}, once more
  \item \emph{Say it:} \emph{acchanŭ aṭutta āḻca oru śastrakriya uṇṭŭ.} — Father has an operation next week.
  \item \emph{Recall:} say the Malayalam for a symptom, then the Malayalam for an infection, then say \emph{śastrakriya} again
\end{itemize}

\begin{tcolorbox}[breakable,colback=teal!4,colframe=teal!35!black,title={Before you move on}]
What does \ml{ശസ്ത്രക്രിയ} mean? (\textquotedblleft{}Surgery; an operation\textquotedblright{}.) Say it once more.
\end{tcolorbox}
\section[\texorpdfstring{\ml{രോഗനിർണ്ണയം} (rōganirṇṇayaṁ)}{രോഗനിർണ്ണയം (rōganirṇṇayaṁ)}]{\ml{രോഗനിർണ്ണയം} (rōganirṇṇayaṁ) — a diagnosis}

\label{lesson:ML-C532-roganirnnayam}



\begin{quote}
Before the new one: say the Malayalam for an infection, then the Malayalam for surgery.
\end{quote}

\subsection*{You'll want to know: \ml{രോഗനിർണ്ണയം}}
\textbf{\ml{രോഗനിർണ്ണയം}} — \emph{rōganirṇṇayaṁ} — \textquotedblleft{}a diagnosis\textquotedblright{}.

\begin{quote}
\ml{പരിശോധനയ്ക്ക്} \ml{ശേഷം} \ml{വൈദ്യൻ} \ml{രോഗനിർണ്ണയം} \ml{നടത്തി}.
\end{quote}

\begin{quote}
\emph{pariśōdhanaykkŭ śēṣaṁ vaidyan rōganirṇṇayaṁ naṭatti.}
\end{quote}

\begin{quote}
After the tests the doctor made a diagnosis.
\end{quote}

\ml{രോഗം} an illness + \ml{നിർണ്ണയം} a settling, a decision.

\begin{cousinweb}
\textbf{\ml{രോഗനിർണ്ണയം}} (\emph{rōganirṇṇayaṁ}) is a learned compound built from Sanskrit \emph{roga}, \textquotedblleft{}disease\textquotedblright{}, and \emph{nirṇaya}, \textquotedblleft{}taking off, settling, a decision\textquotedblright{} (Monier-Williams).
\end{cousinweb}

\begin{grammarlens}[title={What you've built}]
One more word for describing what you lived and felt.
\end{grammarlens}

\subsection*{Your turn}
\begin{itemize}
\raggedright
  \item \emph{Say it:} \emph{rōganirṇṇayaṁ}
  \item \emph{Say it:} \emph{rōganirṇṇayaṁ}, once more
  \item \emph{Say it:} \emph{pariśōdhanaykkŭ śēṣaṁ vaidyan rōganirṇṇayaṁ naṭatti.} — After the tests the doctor made a diagnosis.
  \item \emph{Recall:} say the Malayalam for an infection, then the Malayalam for surgery, then say \emph{rōganirṇṇayaṁ} again
\end{itemize}

\begin{tcolorbox}[breakable,colback=teal!4,colframe=teal!35!black,title={Before you move on}]
What does \ml{രോഗനിർണ്ണയം} mean? (\textquotedblleft{}A diagnosis\textquotedblright{}.) Say it once more.
\end{tcolorbox}
\section[\texorpdfstring{\ml{ചികിത്സിക്കുക} (cikitsikkuka)}{ചികിത്സിക്കുക (cikitsikkuka)}]{\ml{ചികിത്സിക്കുക} (cikitsikkuka) — to treat (a patient)}

\label{lesson:ML-C532-cikitsikkuka}



\begin{quote}
Before the new one: say the Malayalam for surgery, then the Malayalam for a diagnosis.
\end{quote}

\subsection*{You'll want to know: \ml{ചികിത്സിക്കുക}}
\textbf{\ml{ചികിത്സിക്കുക}} — \emph{cikitsikkuka} — \textquotedblleft{}to treat (a patient)\textquotedblright{}.

\begin{quote}
\ml{ആ} \ml{വൈദ്യൻ} \ml{ഒരുപാട്} \ml{രോഗികളെ} \ml{ചികിത്സിച്ചു}.
\end{quote}

\begin{quote}
\emph{ā vaidyan orupāṭŭ rōgikaḷe cikitsiccu.}
\end{quote}

\begin{quote}
That doctor treated a great many patients.
\end{quote}

Past: \ml{ചികിത്സിച്ചു}, \emph{cikitsiccu}.

The noun \ml{ചികിത്സ}, treatment, is already familiar; the patient takes -\ml{എ}.

\begin{cousinweb}
\textbf{\ml{ചികിത്സിക്കുക}} (\emph{cikitsikkuka}) is built from Sanskrit \emph{cikitsā}, \textquotedblleft{}medical attendance, the science of healing\textquotedblright{} (Monier-Williams), with the Malayalam verb ending -\ml{ിക്കുക}.
\end{cousinweb}

\begin{grammarlens}[title={What you've built}]
One more word for describing what you lived and felt.
\end{grammarlens}

\subsection*{Your turn}
\begin{itemize}
\raggedright
  \item \emph{Say it:} \emph{cikitsikkuka}
  \item \emph{Say it:} \emph{cikitsikkuka}, once more
  \item \emph{Say it:} \emph{ā vaidyan orupāṭŭ rōgikaḷe cikitsiccu.} — That doctor treated a great many patients.
  \item \emph{Recall:} say the Malayalam for surgery, then the Malayalam for a diagnosis, then say \emph{cikitsikkuka} again
\end{itemize}

\begin{tcolorbox}[breakable,colback=teal!4,colframe=teal!35!black,title={Before you move on}]
What does \ml{ചികിത്സിക്കുക} mean? (\textquotedblleft{}To treat (a patient)\textquotedblright{}.) Say it once more.
\end{tcolorbox}
\section[(dialogue)]{Words from chapters 525–532 — a review}

\label{lesson:ML-R532-review-525-532}



\begin{quote}
Say the Malayalam for to groan and to burst into tears.
\end{quote}

\subsection*{Your turn}
Say these aloud:

\begin{itemize}
\raggedright
  \item to groan, to burst into tears, a hurt feeling, to burst out laughing, confusion
  \item to be stumped, to get agitated, to thrill, to be amused, to grow weary
  \item kind-hearted, patient by nature, a miser, generous, a loyal person
  \item stubbornness, a hard worker, optimism, sincerity, discipline
  \item good sense, self-assurance, guileless, a bad habit, tolerance
  \item energetic, briskness, obstinacy, a coward, to feel shy
  \item to mock, to boast, to imitate, showing off, to bother
  \item a symptom, an infection, surgery, a diagnosis, to treat
\end{itemize}

\begin{tcolorbox}[breakable,colback=teal!4,colframe=teal!35!black,title={Before you move on}]
To groan? (\textbf{\ml{ഞരങ്ങുക}}, \emph{ñaraṅṅuka}.) To burst into tears? (\textbf{\ml{പൊട്ടിക്കരയുക}}, \emph{poṭṭikkarayuka}.) A hurt feeling? (\textbf{\ml{പരിഭവം}}, \emph{paribhavaṁ}.) To burst out laughing? (\textbf{\ml{പൊട്ടിച്ചിരിക്കുക}}, \emph{poṭṭiccirikkuka}.) Confusion? (\textbf{\ml{ആശയക്കുഴപ്പം}}, \emph{āśayakkuḻappaṁ}.) To be stumped? (\textbf{\ml{കുഴങ്ങുക}}, \emph{kuḻaṅṅuka}.) To get agitated? (\textbf{\ml{ക്ഷോഭിക്കുക}}, \emph{kṣōbhikkuka}.) To thrill? (\textbf{\ml{കോരിത്തരിക്കുക}}, \emph{kōrittarikkuka}.) To be amused? (\textbf{\ml{രസിക്കുക}}, \emph{rasikkuka}.) To grow weary? (\textbf{\ml{മുഷിയുക}}, \emph{muṣiyuka}.) Kind-hearted? (\textbf{\ml{ദയയുള്ള}}, \emph{dayayuḷḷa}.) Patient by nature? (\textbf{\ml{ക്ഷമയുള്ള}}, \emph{kṣamayuḷḷa}.) A miser? (\textbf{\ml{പിശുക്കൻ}}, \emph{piśukkan}.) Generous? (\textbf{\ml{ഉദാരമായ}}, \emph{udāramāya}.) A loyal person? (\textbf{\ml{വിശ്വസ്തൻ}}, \emph{viśvastan}.) Stubbornness? (\textbf{\ml{ശാഠ്യം}}, \emph{śāṭhyaṁ}.) A hard worker? (\textbf{\ml{കഠിനാധ്വാനി}}, \emph{kaṭhinādhvāni}.) Optimism? (\textbf{\ml{ശുഭാപ്തിവിശ്വാസം}}, \emph{śubhāptiviśvāsaṁ}.) Sincerity? (\textbf{\ml{ആത്മാർത്ഥത}}, \emph{ātmārtthata}.) Discipline? (\textbf{\ml{അച്ചടക്കം}}, \emph{accaṭakkaṁ}.) Good sense? (\textbf{\ml{വിവേകം}}, \emph{vivēkaṁ}.) Self-assurance? (\textbf{\ml{തന്റേടം}}, \emph{tanṟēṭaṁ}.) Guileless? (\textbf{\ml{നിഷ്കളങ്കമായ}}, \emph{niṣkaḷaṅkamāya}.) A bad habit? (\textbf{\ml{ദുശ്ശീലം}}, \emph{duśśīlaṁ}.) Tolerance? (\textbf{\ml{സഹിഷ്ണുത}}, \emph{sahiṣṇuta}.) Energetic? (\textbf{\ml{ഊർജ്ജസ്വലമായ}}, \emph{ūrjjasvalamāya}.) Briskness? (\textbf{\ml{ചുറുചുറുക്ക്}}, \emph{cuṟucuṟukkŭ}.) Obstinacy? (\textbf{\ml{വാശി}}, \emph{vāśi}.) A coward? (\textbf{\ml{ഭീരു}}, \emph{bhīru}.) To feel shy? (\textbf{\ml{നാണിക്കുക}}, \emph{nāṇikkuka}.) To mock? (\textbf{\ml{പരിഹസിക്കുക}}, \emph{parihasikkuka}.) To boast? (\textbf{\ml{വീമ്പിളക്കുക}}, \emph{vīmpiḷakkuka}.) To imitate? (\textbf{\ml{അനുകരിക്കുക}}, \emph{anukarikkuka}.) Showing off? (\textbf{\ml{പൊങ്ങച്ചം}}, \emph{poṅṅaccaṁ}.) To bother? (\textbf{\ml{ശല്യപ്പെടുത്തുക}}, \emph{śalyappeṭuttuka}.) A symptom? (\textbf{\ml{ലക്ഷണം}}, \emph{lakṣaṇaṁ}.) An infection? (\textbf{\ml{അണുബാധ}}, \emph{aṇubādha}.) Surgery? (\textbf{\ml{ശസ്ത്രക്രിയ}}, \emph{śastrakriya}.) A diagnosis? (\textbf{\ml{രോഗനിർണ്ണയം}}, \emph{rōganirṇṇayaṁ}.) To treat? (\textbf{\ml{ചികിത്സിക്കുക}}, \emph{cikitsikkuka}.)
\end{tcolorbox}
